% What does PM reveal? What do we now know about complex systems that we didn't before?
\section{Conclusion}
\label{sec:conclusion}
% Reveal: It is what sits beneath existing disciplines
Fluctuation theorems, master equations, and Langevin dynamics describe the statistics of transitions; the architecture of persistence describes the structure that must exist before any such statistics can be defined. For the fields of adaptation and self-organization, this shifts the fundamental question: before a system can adapt or self-organize, it must possess a minimal architecture permitted to exist.

% What we showed (grounds the shift)
We have developed that architecture and a mechanism of selection under the stated dynamical assumptions. The six pillars organize the requirements of persistence; the Operational Action measures cumulative dissipation; and the Selection Principle describes the resulting survival bias when failure hazard is proportional to that action.

% Reveal: Transfer
The shared architecture provides a starting point for transferring knowledge across disciplines. Because this universal mapping projects every persistent system onto the same six-pillar substrate, it functions as a Rosetta stone across disciplines. Hard-won solutions in one domain become candidates for transfer to another: protocol engineering that prevents signal decay in a fiber-optic network can suggest mechanisms for preventing fragmentation in a corporate hierarchy. In emerging domains such as AI, where the pillars are present but unnamed, this vocabulary potentially lets systems be engineered for resilience without re-deriving structural knowledge from scratch.

% Reveal: Compare
The same mapping enables direct quantitative comparison among competing systems in a shared environment. By evaluating how competing organisms, firms, or technical architectures manage their finite dissipation budgets, relative operational resilience can be assessed prior to empirical failure.

% Reveal: Quantify (failure vocabulary becomes shared and rigorous)
Organizational decline, ecosystem collapse, and the limits of biological longevity have been described through domain-specific frameworks. By defining exactly what is required to persist, this mechanics recasts them as structural failures under a single, unified mechanics: disparate phenomena become comparable instances of the same dissipation constraint, rather than loose analogies.

% Reveal: Predict
The Complexity Trade-off is a pre-observational constraint on structural complexity, coordination speed, and energetic efficiency. At a binding dissipation ceiling, increasing one factor requires a compensating reduction in the product of the other two; given the ceiling, any two positive factors bound the third.

% Future Payoff
The most consequential limiting case is fundamental physics: whether standard physical law is itself a substrate-specific instantiation of this same invariant. Further exploration will determine whether this architectural boundary holds universally. Candidate mathematical languages for testing this include Information Geometry (continuous substrates), Arithmetic Geometry (discrete substrates), and Sheaf Theory (compositional substrates).

% Future work
Natural avenues for the extension include systems composed of nested persistent subsystems, what it takes for a system to arise from disorder, and a fuller account of decay itself: this paper identifies failure modes, without characterizing the trajectory a system follows on the way there, or whether decay has its own structure worth deriving.

% The mechanics of persistence provides a missing piece to science
Persistence Mechanics does not explain why structure arises. It explains why only certain architectures remain to be observed. Persistence requires no intent, no optimization, and no forward-looking agency, only that the substrate continuously filter non-viable configurations. What remains observable is the minimal architecture of persistence.